\chapter{Điểm số, sĩ diện và áp lực thật}
\markboth{Điểm số, sĩ diện và áp lực thật}{Điểm số, sĩ diện và áp lực thật}

\section{Điểm thấp là do đề gài em thôi}
\begin{truyen}{Điểm thấp là do đề gài em thôi}{Classroom Talk}
\chuhoa{H}{ọc sinh: ``Đề này gài quá trời, chứ không phải em dốt.''}
\textit{Student: ``This test was full of traps. It is not that I am bad at the subject.''}\\[4pt]

\teacherpair{Đề có gài thật hay không, mình xem lại sau. Nhưng em né một chuyện hơi rõ: phần em mất điểm toàn nằm chỗ em hổng.}{Whether the test was tricky or not, we can check that later. But you are avoiding one obvious fact: every place you lost points is exactly where your knowledge is weak.}

\studentpair{Nhưng em nhìn vô là em biết hết mà.}{But when I looked at it, I felt like I knew everything.}

\teacherpair{Biết cảm giác khác với biết chắc. Điểm thấp nhiều khi đau vì nó đụng đúng chỗ mình tưởng mình ổn.}{Feeling like you know something is different from truly knowing it. Low grades often hurt because they hit the exact place where we thought we were solid.}
\end{truyen}

\begin{baihoc}
Đổ cho đề là cách giữ sĩ diện nhanh nhất. Nhưng nhìn thẳng vào lỗ hổng mới là cách tăng điểm thật.
\textit{(Blaming the test protects pride quickly. Facing the gap improves results for real.)}
\end{baihoc}

\begin{ghinhoanh}
\textbf{Short English to remember:}

Blaming the test saves pride.

Facing the gap builds skill.
\end{ghinhoanh}

\begin{tuhocnhanh}
1. Em đang giận đề hay giận chính chỗ mình hổng? \textit{(Are you angry at the test, or at your own weak spot?)}

2. Em đã xem kỹ mình sai ở dạng nào chưa? \textit{(Have you checked exactly what type of mistake you made?)}
\end{tuhocnhanh}
\ngancach

\section{Em hiểu bài mà vô làm lại không ra}
\begin{truyen}{Em hiểu bài mà vô làm lại không ra}{Classroom Talk}
\chuhoa{H}{ọc sinh: ``Lúc thầy giảng em hiểu lắm. Tới khi làm bài em đứng hình. Vậy lỗi do em run chứ đâu phải em không biết.''}
\textit{Student: ``When you explain it, I understand clearly. But when I try to do it myself, I freeze. So it is just nerves, not lack of understanding, right?''}\\[4pt]

\teacherpair{Nếu biết thật, em có cần phụ thuộc vào cảm giác hiểu trong lúc nghe không?}{If you truly knew it, would you still have to depend on the feeling of understanding while listening?}

\studentpair{Chắc em hiểu kiểu tưởng là hiểu.}{Maybe I only understand it in the way people think they understand something.}

\teacherpair{Đúng. Nhiều bạn nhầm giữa nghe thấy hợp lý với tự làm được bằng tay mình.}{Exactly. Many students confuse hearing something that sounds reasonable with actually being able to do it by themselves.}
\end{truyen}

\begin{baihoc}
Hiểu thật là khi ngồi một mình vẫn làm được, không có thầy nhắc từng bước.
\textit{(Real understanding means being able to do it alone, without step-by-step prompting.)}
\end{baihoc}

\begin{ghinhoanh}
\textbf{Short English to remember:}

Feeling clear is not the same as being able to do it.

Practice exposes fake understanding.
\end{ghinhoanh}

\begin{tuhocnhanh}
1. Em đang hiểu bằng tai hay hiểu bằng tay? \textit{(Do you understand with your ears, or with your hands?)}

2. Sau giờ học, em có tự làm lại một bài tương tự không? \textit{(After class, do you solve a similar problem on your own?)}
\end{tuhocnhanh}
\ngancach

\section{Bạn kia học ít hơn mà điểm vẫn cao hơn}
\begin{truyen}{Bạn kia học ít hơn mà điểm vẫn cao hơn}{Classroom Talk}
\chuhoa{H}{ọc sinh: ``Em thấy bất công ghê. Bạn kia đăng story suốt mà điểm vẫn cao. Em ngồi hoài mà điểm không bằng.''}
\textit{Student: ``It feels so unfair. That person posts stories all the time and still gets high scores. I sit and study for ages and still cannot catch up.''}\\[4pt]

\teacherpair{Em thấy thời gian nó khoe, hay em thấy hết thời gian nó học?}{Are you seeing all the time that person studies, or only the time they choose to show off?}

\studentpair{Thì... em đâu ở cạnh bạn cả ngày.}{Well... I am obviously not beside that person all day.}

\teacherpair{Mình hay so cái hậu trường tệ nhất của mình với cái phần thiên hạ cho mình thấy. So kiểu đó chỉ có tự làm mình tụt tinh thần.}{We often compare the messiest backstage part of our own lives with the polished part other people choose to show. That kind of comparison only drags your spirit down.}
\end{truyen}

\begin{baihoc}
So sánh thiếu dữ liệu chỉ làm mình cay, chứ không làm mình khá lên.
\textit{(Comparison without full data only breeds frustration, not improvement.)}
\end{baihoc}

\begin{ghinhoanh}
\textbf{Short English to remember:}

You rarely see the whole story.

Comparison without context is a trap.
\end{ghinhoanh}

\begin{tuhocnhanh}
1. Em đang so với thực tế hay so với ảo giác mình tự dựng? \textit{(Are you comparing with reality, or with an illusion?)}

2. Điều gì trong cách học của em cần sửa cụ thể? \textit{(What specific part of your study method needs fixing?)}
\end{tuhocnhanh}
\ngancach

\section{Bị điểm thấp, em sợ nhất là về nhà}
\begin{truyen}{Bị điểm thấp, em sợ nhất là về nhà}{Classroom Talk}
\chuhoa{H}{ọc sinh: ``Em không sợ thầy bằng sợ đem bài này về nhà. Ở lớp em bị quê một lần, về nhà em bị nhắc cả tuần.''}
\textit{Student: ``I am less afraid of you than of taking this paper home. In class I am embarrassed once, but at home I will hear about it all week.''}\\[4pt]

\teacherpair{Thầy hiểu. Nhưng trốn xem điểm không làm bài tốt lên. Mình cần nghĩ cách sửa bài và cách báo với nhà cho đàng hoàng.}{I understand that. But hiding from the grade will not improve the result. We need to think about how to repair the mistake and how to tell your family properly.}

\studentpair{Báo sao cho đỡ nổ?}{How do I tell them without everything exploding?}

\teacherpair{Đừng chỉ đưa điểm. Đưa luôn kế hoạch sửa. Người lớn bớt giận khi thấy em chưa buông.}{Do not just show the score. Show the plan to improve it too. Adults calm down more easily when they see you have not given up.}
\end{truyen}

\begin{baihoc}
Điểm thấp không đáng sợ bằng thái độ trốn tránh sau điểm thấp.
\textit{(A low grade is less damaging than the avoidance that follows it.)}
\end{baihoc}

\begin{ghinhoanh}
\textbf{Short English to remember:}

Bring the result.

Bring a repair plan too.
\end{ghinhoanh}

\begin{tuhocnhanh}
1. Khi bị điểm thấp, em thường trốn hay đối diện? \textit{(When you get a low grade, do you hide or face it?)}

2. Nếu phải nói với gia đình, kế hoạch sửa của em là gì? \textit{(If you had to tell your family, what would your repair plan be?)}
\end{tuhocnhanh}
\ngancach